\section{Letters, brackets, and the equals sign}\label{ch01:sec:letters}
A letter such as $n$ that stands for an integer whose value we have not fixed is called a \emph{variable}. A number written directly in front of a letter means multiplication: $2n$ means $2\times n$. Two letters written side by side are multiplied too, so $ab$ means $a\times b$. A raised dot also means multiplication: $2\cdot4$ is $2\times4$. Between two numerals the dot cannot be left out, because $24$ already means twenty-four.

The quantities multiplied in a product are its \emph{factors}, and the quantities added in a sum are its \emph{terms}: the product $2n$ has the factors $2$ and $n$, and the sum $2n+1$ has the terms $2n$ and $1$. The square $n^2$ means $n\times n$. Parentheses, also called brackets, group an expression: $2(n+1)$ is twice the whole quantity $n+1$. It equals $2n+2$, not $2n+1$, because the $1$ inside the brackets is doubled too. At $n=5$, for instance, $2(5+1)=12$, while $2\cdot5+1=11$.

An equals sign says that the expressions on either side of it have the same value; it does not mean ``and now calculate the next thing.'' The chain $2(3+1)=2\cdot4=8$ is correct, because all three expressions have the value $8$. The chain $2+3=5\cdot2=10$ is wrong, because $2+3$ is not equal to $5\cdot2$. The writer meant two steps, and should have written two equations.

We will use the familiar rules of arithmetic freely. Addition and multiplication can each be done in either order and grouped in either way, and multiplication distributes over addition:
\begin{gather*}
a+b=b+a,\qquad (a+b)+c=a+(b+c),\\
ab=ba,\qquad (ab)c=a(bc),\qquad a(b+c)=ab+ac.
\end{gather*}
The last rule is called the \emph{distributive rule}. Because multiplication can be done in either order, it also works with the sum written first: $(a+b)c=ac+bc$. These rules hold for all integers, and more generally for all the real numbers described in Section~\ref{ch01:sec:notation}.

A fraction $a/b$ means $a$ divided by $b$, and it is defined only when $b\ne0$ (the symbol $\ne$ means ``is not equal to''). If two expressions are equal, then adding the same number to both, subtracting it from both, or multiplying both by it keeps them equal, and so does dividing both by the same nonzero number.

\pauseq{An assistant starts from the equation $x^2=x$, divides both sides by $x$, and concludes that $x=1$. Has it found every solution?}{No. Dividing by $x$ is allowed only when $x\ne0$, so the argument says nothing about $x=0$; forgetting such a case is a common way to lose solutions. A complete argument treats two cases. If $x=0$, the equation holds, since $0^2=0$. If $x\ne0$, dividing both sides by $x$ gives $x=1$, and indeed $1^2=1$. So the solutions are exactly $0$ and $1$.}
